\documentclass[11pt]{article}
\usepackage[margin=1in]{geometry}
\usepackage{amsmath,amssymb,amsthm}
\usepackage{booktabs}
\usepackage{hyperref}

\theoremstyle{plain}
\newtheorem{thm}{Theorem}
\newtheorem{lem}[thm]{Lemma}
\newtheorem{cor}[thm]{Corollary}
\newtheorem{prop}[thm]{Proposition}
\theoremstyle{definition}
\newtheorem{defn}[thm]{Definition}
\theoremstyle{remark}
\newtheorem{rem}[thm]{Remark}
\newtheorem{conj}[thm]{Conjecture}
\theoremstyle{plain}
\newtheorem{gapthm}{Gap}
\newtheorem{retr}{Retraction}

\newcommand{\bI}{\mathbf{I}}
\newcommand{\bM}{\mathbf{M}}
\newcommand{\bw}{\mathbf{w}}
\newcommand{\At}{\widetilde{A}}
\newcommand{\Ah}{\widehat{A}}
\newcommand{\ord}{\operatorname{ord}}
\newcommand{\Res}{\operatorname{Res}}
\newcommand{\sm}{\setminus}

\title{The long-root multiplier is a coboundary,\\[2mm]
\large and therefore the long-root poles of the structure constants are
\emph{removable}:\\
a retraction of my own pre-registered answer, and a confirmation of
Rick's hunch\\[2mm]
\normalsize
with the exact family of renormalisations that remove them ---
$\chi(x)\chi(1/x)=\rho^{-1}$ \emph{and} $\chi(\varepsilon)=1$ for $\varepsilon^h=1$}
\author{Clio}
\date{17 October 2026}

\makeatletter
%% protocol-header.tex -- PROTOCOL §2.3 page-1 block, as a TEMPLATE not an instance.
%%
%% PROTOCOL §2.3: "Every PDF must state, on its first page: author, date,
%% recipient, a title, and the commit hash of the work-in-progress push it
%% corresponds to.  A PDF with no commit hash cannot be checked against the
%% record and should be treated as a draft opinion, not a result."
%%
%% Written 2026-10-16 because the fault had recurred twice and I had twice
%% fixed only the instance:
%%   - 2026-10-14: the Q473 PDF went to Rick with no recipient and no hash.
%%   - 2026-10-14: the Q460 PDF went to Rick saying "Recipient: Robin Langer".
%%     He noticed, and took the commit hash as authority instead.
%% Both times I repaired the one file and wrote "fix the template, not just the
%% instance" in my notes.  This file is that fix.  \input it; do not retype it.
%%
%% USAGE, immediately after \maketitle:
%%   %% protocol-header.tex -- PROTOCOL §2.3 page-1 block, as a TEMPLATE not an instance.
%%
%% PROTOCOL §2.3: "Every PDF must state, on its first page: author, date,
%% recipient, a title, and the commit hash of the work-in-progress push it
%% corresponds to.  A PDF with no commit hash cannot be checked against the
%% record and should be treated as a draft opinion, not a result."
%%
%% Written 2026-10-16 because the fault had recurred twice and I had twice
%% fixed only the instance:
%%   - 2026-10-14: the Q473 PDF went to Rick with no recipient and no hash.
%%   - 2026-10-14: the Q460 PDF went to Rick saying "Recipient: Robin Langer".
%%     He noticed, and took the commit hash as authority instead.
%% Both times I repaired the one file and wrote "fix the template, not just the
%% instance" in my notes.  This file is that fix.  \input it; do not retype it.
%%
%% USAGE, immediately after \maketitle:
%%   %% protocol-header.tex -- PROTOCOL §2.3 page-1 block, as a TEMPLATE not an instance.
%%
%% PROTOCOL §2.3: "Every PDF must state, on its first page: author, date,
%% recipient, a title, and the commit hash of the work-in-progress push it
%% corresponds to.  A PDF with no commit hash cannot be checked against the
%% record and should be treated as a draft opinion, not a result."
%%
%% Written 2026-10-16 because the fault had recurred twice and I had twice
%% fixed only the instance:
%%   - 2026-10-14: the Q473 PDF went to Rick with no recipient and no hash.
%%   - 2026-10-14: the Q460 PDF went to Rick saying "Recipient: Robin Langer".
%%     He noticed, and took the commit hash as authority instead.
%% Both times I repaired the one file and wrote "fix the template, not just the
%% instance" in my notes.  This file is that fix.  \input it; do not retype it.
%%
%% USAGE, immediately after \maketitle:
%%   \input{protocol-header}
%%   \protocolheader{Rick (grandpa-rick)}{clio-vega/proofs}{60e4ce3}
%%
%% The recipient argument is MANDATORY and has no default on purpose: a default
%% is what produced "Recipient: Robin Langer" on a PDF addressed to Rick.

%% \maketitle CLEARS \@date (and \@author, \@title) in the article class, so a
%% macro called AFTER \maketitle that reads \@date prints an EMPTY row.  Caught
%% 2026-10-16 only because I ran pdftotext on page 1 instead of trusting rc=0:
%% the build was clean and the Date field was blank.  Hence: snapshot the date
%% HERE, at \input time in the preamble, where \@date is still live.
\let\protocol@date\@date

\newcommand{\protocolheader}[3]{%
  \begin{center}
  \fbox{\begin{minipage}{0.92\textwidth}
  \small
  \textbf{PROTOCOL §2.3 record.}\\[1mm]
  \begin{tabular}{@{}ll@{}}
  Author     & Clio \\
  Date       & \protocol@date \\
  Recipient  & #1 \\
  Repository & \texttt{#2} \\
  Commit     & \texttt{#3} \\
  \end{tabular}\\[1mm]
  This PDF corresponds to the pushed commit above.  Where this text and the
  commit disagree, \textbf{the commit is authoritative}.
  \end{minipage}}
  \end{center}
  \vspace{2mm}%
}

%%   \protocolheader{Rick (grandpa-rick)}{clio-vega/proofs}{60e4ce3}
%%
%% The recipient argument is MANDATORY and has no default on purpose: a default
%% is what produced "Recipient: Robin Langer" on a PDF addressed to Rick.

%% \maketitle CLEARS \@date (and \@author, \@title) in the article class, so a
%% macro called AFTER \maketitle that reads \@date prints an EMPTY row.  Caught
%% 2026-10-16 only because I ran pdftotext on page 1 instead of trusting rc=0:
%% the build was clean and the Date field was blank.  Hence: snapshot the date
%% HERE, at \input time in the preamble, where \@date is still live.
\let\protocol@date\@date

\newcommand{\protocolheader}[3]{%
  \begin{center}
  \fbox{\begin{minipage}{0.92\textwidth}
  \small
  \textbf{PROTOCOL §2.3 record.}\\[1mm]
  \begin{tabular}{@{}ll@{}}
  Author     & Clio \\
  Date       & \protocol@date \\
  Recipient  & #1 \\
  Repository & \texttt{#2} \\
  Commit     & \texttt{#3} \\
  \end{tabular}\\[1mm]
  This PDF corresponds to the pushed commit above.  Where this text and the
  commit disagree, \textbf{the commit is authoritative}.
  \end{minipage}}
  \end{center}
  \vspace{2mm}%
}

%%   \protocolheader{Rick (grandpa-rick)}{clio-vega/proofs}{60e4ce3}
%%
%% The recipient argument is MANDATORY and has no default on purpose: a default
%% is what produced "Recipient: Robin Langer" on a PDF addressed to Rick.

%% \maketitle CLEARS \@date (and \@author, \@title) in the article class, so a
%% macro called AFTER \maketitle that reads \@date prints an EMPTY row.  Caught
%% 2026-10-16 only because I ran pdftotext on page 1 instead of trusting rc=0:
%% the build was clean and the Date field was blank.  Hence: snapshot the date
%% HERE, at \input time in the preamble, where \@date is still live.
\let\protocol@date\@date

\newcommand{\protocolheader}[3]{%
  \begin{center}
  \fbox{\begin{minipage}{0.92\textwidth}
  \small
  \textbf{PROTOCOL §2.3 record.}\\[1mm]
  \begin{tabular}{@{}ll@{}}
  Author     & Clio \\
  Date       & \protocol@date \\
  Recipient  & #1 \\
  Repository & \texttt{#2} \\
  Commit     & \texttt{#3} \\
  \end{tabular}\\[1mm]
  This PDF corresponds to the pushed commit above.  Where this text and the
  commit disagree, \textbf{the commit is authoritative}.
  \end{minipage}}
  \end{center}
  \vspace{2mm}%
}

\makeatother

\begin{document}
\maketitle
\protocolheader{Rick (\texttt{grandparick20@gmail.com})}{clio-vega/proofs}{COMMITHASH}

\begin{abstract}
On 2026-10-17 Rick (UID 819) proved \textbf{Lemma QC}: on the long-root divisor
$H_{ab}=\{w_aw_b=1\}$, Motegi's normalised type-$C$ weight functions satisfy
$\At_{\bM,\bI}=g\cdot\At_{\bM\sm\{a,b\},\bI}$ for \emph{every} $\bI$, with
$g=(1-q)^2/((1-qw_a^2)(1-qw_a^{-2}))$.  This is my $\rho^C$, reached by a
different route, and it upgrades my Lemma A from a divisibility to an identity.
He attached a hunch: \emph{``Clio's poles measure the normalisation defect; a
geometric normalisation would have no long-root poles at all.''}

I pre-registered the answer \textbf{no}, with a mechanism.  \textbf{The
pre-registered answer is wrong and is retracted here.}  Rick is right, and the
reason is a one-line observation about $\rho$ that I had already proved and had
not read correctly: writing $\varphi(v):=(1-q)/(1-qv^{h})$,
\[
   \rho^X=\varphi(w_a)\,\varphi(w_b),
\]
a product of a \emph{single universal one-variable function} evaluated at the two
colliding indices.  That is exactly the statement that the collapse multiplier is
a \textbf{coboundary}.  Setting $\psi:=\varphi^{-1}$ and
$\Psi(\bM):=\prod_{m\in\bM}\psi(w_m)$, the renormalised basis
$\Ah_{\bM,\bI}:=\At_{\bM,\bI}\Psi(\bM)$ collapses with multiplier~$1$ for every
long-root pair simultaneously (Theorem~\ref{thm:A}), \emph{and} every one of its
structure constants is regular on every long-root divisor
(Theorem~\ref{thm:B}).  Measured at $N=3$, type $C$: the baseline census returns
$13$ non-Laurent constants out of $115$, with poles on long-root divisors only;
after the renormalisation the count is $\mathbf{0/115}$.

The exact family is determined (Theorem~\ref{thm:C}): among renormalisations
$\Psi_\chi(\bM)=\prod_{m\in\bM}\chi(w_m)$, the long-root poles vanish
\emph{iff} $\chi(x)\chi(1/x)=\rho^{-1}|_H$, and --- a condition the session did
not anticipate and the data forced --- the \emph{short}-root divisors
$\{w_m=\varepsilon\}$, $\varepsilon^{h}=1$, stay pole-free \emph{iff}
$\chi(\varepsilon)=1$.  The short-root collapse multiplier is measured to be
exactly~$1$.  Both conditions are met by $\psi$ for the same reason:
$\psi(\varepsilon)=1$ precisely when $\varepsilon^h=1$, which is the locus
$\{\rho=1\}$.

Consequences.  (i) Q491's criterion ``$c^\nu_{\lambda\mu}$ Laurent $\iff$
$|\nu|-\max(|\lambda|,|\mu|)\le1$'' is a property of the
$\mathcal{E}^C$-normalised basis, \textbf{not} of the algebra: under $\psi$ it
becomes vacuous.  The negative half of Q500 --- a ``Pfaffian-category
obstruction'' read off the non-Laurentness --- is withdrawn in the form in which
I stated it.  (ii) What survives, and is the right invariant, is \emph{not} the
pole but the \emph{class}: the poles are removable because $\rho$ is a
coboundary, and had $\rho$ not split as $\varphi(w_a)\varphi(w_b)$ no
point-diagonal renormalisation could have removed them.  The obstruction I had
been measuring was a cochain.  The class is zero.
\end{abstract}

\tableofcontents

%=====================================================================
\section{What this session set out to test, and what happened}

The brief (\texttt{state/PROVE.md}, 2026-10-17 c2) set three targets.
Theorem~A, expected easy: $\rho^X$ is a coboundary.  Theorem~B,
\emph{the actual question}: does the same renormalisation remove the poles of
the \emph{structure constants}?  Theorem~C, conditional on B: is the pole order
invariant across the whole family of diagonal renormalisations?

The brief pre-registered the answer to B as \textbf{no}, with this mechanism,
quoted verbatim:

\begin{quote}\small
``the renormalisation is a diagonal rescaling of the basis by
$\prod_{m\in\bM}\psi(w_m)$, so a structure constant transforms by a ratio of
$\psi$-products.  $\psi(w)=(1-qw^h)/(1-q)$ has \dots\ no zeros or poles on the
collision divisor $w_aw_b=1$ at generic $w_a$ \dots\ A factor that is regular and
non-vanishing at the generic point of $H_{ab}$ \textbf{cannot cancel a pole
there}.''
\end{quote}

The final sentence is true.  The sentence before it is false, and
\S\ref{sec:error} isolates the error exactly: the product is indexed by
$m\in\bM$ --- the \emph{specialisation set}, i.e.\ the \emph{point} --- and a
rescaling indexed by the point is not a rescaling of the basis.  A rescaling of
the basis is indexed by $\bI$, the \emph{label}.  The two operations do different
things to structure constants, and, decisively, \textbf{Theorem~A is true only
for the point-indexed reading while the Theorem~B argument is valid only for the
label-indexed one}.  One symbol, two readings, two paragraphs apart, each
locally plausible.

So Theorem~B is false as pre-registered, and its negation is a theorem.  Theorem~C
is then settled in the opposite direction to the brief's guess: the pole order is
\emph{not} invariant across the family --- it is $1$ off a distinguished coset of
renormalisations and $0$ on it --- and the coset is computed exactly.

%=====================================================================
\section{Task 0: which object is Rick's \texorpdfstring{$\At$}{Atilde}?}
\label{sec:task0}

The brief imposed a hard 30-minute gate here, because the 2026-10-17 c1
retraction was caused by exactly this class of error (a type-$B$ weight function
divided by the type-$C$ normaliser), and because Rick's $g$ matching my $\rho^C$
has the shape of a one-coordinate match.  I read
\texttt{peers/rick/proofs/2026-10-17-q488-converse-order-q1-proofs.pdf}
first-hand.  All three answers are ``same object''.

\begin{enumerate}
\item \textbf{Normaliser.}  His fact (M2), \S1: ``\texttt{E\textasciicircum C}
  (L1057--1062); $\At=WC/E^C$.''  Type $C$ weight function divided by the type
  $C$ normaliser.  Mine likewise.  (This is the coordinate on which c1 failed,
  and it is clean here.)
\item \textbf{Index convention.}  His \S1 defines
  $\phi_{\bI}(\bM):=\At_{\bM,\bI}:=\widetilde{WC}_{N,n,|\bI|}(w_{\bM}\mid
  w\mid\bI)$.  The subscript $\bM$ enters as the \emph{specialisation}
  $u=w_{\bM}$ and $\bI$ as the \emph{label} carried by $W$.  That is my
  convention; his UID~819 Ask~(1) asked me to check precisely this and the answer
  is that his $\bI$ is my $\bI$.
\item \textbf{Divisor.}  His $H_{ab}=\{w_aw_b=1\}$, local parameter
  $1-1/(w_aw_b)$; my $H$ is the same divisor with the same parameter.
\end{enumerate}

\noindent
I also checked the arithmetic rather than only the names: with $h=2$,
$\rho^C=(1-q)^2/((1-qw_a^2)(1-qw_b^2))$ restricted to $w_b=w_a^{-1}$ is
$(1-q)^2/((1-qw_a^2)(1-qw_a^{-2}))=g$.  Identical.

\begin{rem}[grading]
Lemma QC is graded \textbf{proved} and not \texttt{peer-claimed}, because I read
his proof (\S3 of the proof file) rather than its statement.  His \S3 is a
term-by-term comparison of the symmetrisation formula at $n=s$ and $n=s-2$: the
head-pair factor $(u_ju_k-1)$ kills every permutation putting both of $x,1/x$ in
the head, the surviving head selections are exactly those of the $(s-2)$-formula,
and the leftover bracket is independent of the head selection.  I have not
re-derived it line by line; I have checked its two consumable outputs
(the multiplier, and uniformity in $\bI$) against my own engine, which reaches
the same $\rho^C$ by a different route.  Independent corroboration in the strong
sense: different author, different derivation.
\end{rem}

%=====================================================================
\section{Setup}

Fix $N\ge1$.  Let $F:=\mathbb{C}(q,w_1,\dots,w_N)$, let
$\mathcal{P}:=2^{[N]}$ be the set of subsets of $[N]$, and let
\[
  \mathcal{A}:=F^{\mathcal{P}}
\]
be the algebra of $F$-valued functions on $\mathcal{P}$ under pointwise
operations --- the ``GKM algebra'' of the $2^N$ fixed points.  For
$\bI\subseteq[N]$ put $\phi_{\bI}\in\mathcal{A}$,
$\phi_{\bI}(\bM):=\At_{\bM,\bI}$, Motegi's normalised weight function
($\At^X:=W^X/\mathcal{E}^X$, evaluated at $u=w_{\bM}$; $n$-independent by
stability, Motegi L2823--2848).  Two facts from Motegi, arXiv:2610.15344, both
deep-read:

\begin{description}
\item[(M3) Triangularity (Thm 2.6(i), L603--620).]
  $\At_{\bM,\bI}=0$ unless $\bM\ge\bI$ in Motegi's order ($|\bM|\ge|\bI|$ and the
  top $|\bI|$ elements of $\bM$ dominate $\bI$ entrywise).
\item[(M4) The diagonal (eq.\ \texttt{normalizationWC}, L1078).]
  $\At_{\bI,\bI}$ is a product over the positive roots of $C_N$.  Its factors are
  $\prod_{j\le k}(1-w_{i_j}^{-1}w_{i_k}^{-1})$ over $\bI$ --- note $j\le k$, so
  the \emph{diagonal} term $(1-w_{i}^{-2})$ occurs --- together with $A$-type
  factors $(1-w_j/w_{i_l})$ and $q$-factors $(1-q\cdot\text{monomial})$.
\end{description}

By (M3)+(M4), $\{\phi_{\bI}\}_{\bI\subseteq[N]}$ is an $F$-basis of
$\mathcal{A}$, and the \emph{structure constants} $c^{\nu}_{\lambda\mu}\in F$ are
defined by $\phi_\lambda\phi_\mu=\sum_\nu c^{\nu}_{\lambda\mu}\phi_\nu$.

Two families of divisors in $w$-space will matter.  For $a<b$ the
\textbf{long-root divisor} $H_{ab}:=\{w_aw_b=1\}$, with local ring
$\mathcal{O}_{ab}\subset F$, uniformiser $t_{ab}:=1-(w_aw_b)^{-1}$, residue field
$k_{ab}$; for $\bM\supseteq\{a,b\}$ write $\bM'':=\bM\sm\{a,b\}$.  For $m\in[N]$
and $\varepsilon^{h}=1$ the \textbf{short-root divisor}
$S_m^{\varepsilon}:=\{w_m=\varepsilon\}$; for $\bM\ni m$ write
$\bM^{-}:=\bM\sm\{m\}$.  On $H_{ab}$ we abbreviate $x:=w_a$, so $w_b=x^{-1}$.

\begin{description}
\item[(QC) Long-root collapse.]  For all $a<b$, all $\bM\supseteq\{a,b\}$, all
  $\bI$: $\ \phi_{\bI}(\bM)\equiv\rho^X\,\phi_{\bI}(\bM'')\pmod{w_aw_b-1}$, where
  \[
     \rho^X=\frac{(1-q)^2}{(1-qw_a^{h})(1-qw_b^{h})},
     \qquad h=2\ (\text{type }C),\quad h=1\ (\text{type }B).
  \]
\end{description}
Type $C$ is Rick's Lemma QC (\S\ref{sec:task0}) and my
\texttt{2026-10-16-c2-longroot-multiplier.tex} Thm \texttt{collapse}; type $B$ is
\texttt{2026-10-17-typeB-multiplier-and-boundary-entry.tex} Thm \texttt{rhoB}.
Write $g:=\rho^X|_{H_{ab}}=(1-q)^2/((1-qx^{h})(1-qx^{-h}))\in k_{ab}^{\times}$.

%=====================================================================
\section{Theorem A: the multiplier is a coboundary}

\begin{thm}[$\rho$ splits]\label{thm:A}
Put $\varphi(v):=\dfrac{1-q}{1-qv^{h}}\in\mathbb{C}(q)(v)$ and
$\psi:=\varphi^{-1}$, i.e.\ $\psi(v)=\dfrac{1-qv^{h}}{1-q}$.  Then
\[
  \rho^X=\varphi(w_a)\,\varphi(w_b)
  \qquad\text{for every pair } a<b,
\]
with the \emph{same} one-variable $\varphi$ for every pair.  Consequently, setting
$\Psi(\bM):=\prod_{m\in\bM}\psi(w_m)$ and
$\Ah_{\bM,\bI}:=\At_{\bM,\bI}\,\Psi(\bM)$, we have for every long-root pair
$a<b$ simultaneously and every $\bI$
\[
  \Ah_{\bM,\bI}\big|_{H_{ab}}=\Ah_{\bM'',\bI}\big|_{H_{ab}}:
\]
honest, untwisted GKM collapse, multiplier exactly $1$.
\end{thm}

\begin{proof}
$\varphi(w_a)\varphi(w_b)=(1-q)^2/((1-qw_a^h)(1-qw_b^h))=\rho^X$ by inspection.
For the consequence, $\Psi(\bM)=\Psi(\bM'')\psi(w_a)\psi(w_b)$ since
$\bM=\bM''\sqcup\{a,b\}$, so by (QC)
\[
  \Ah_{\bM,\bI}\big|_{H}
  =\psi(w_a)\psi(w_b)\,\Psi(\bM'')\cdot\rho^X\,\At_{\bM'',\bI}\big|_{H}
  =\bigl(\varphi(w_a)\varphi(w_b)\bigr)^{-1}\rho^X\cdot\Ah_{\bM'',\bI}\big|_{H}
  =\Ah_{\bM'',\bI}\big|_{H}. \qedhere
\]
\end{proof}

\begin{rem}[what the content is]
The algebra is one line; the content is the \emph{quantifier}.  A priori the
collapse multiplier on $H_{ab}$ is a function attached to the \emph{pair}
$\{a,b\}$.  Theorem~\ref{thm:A} says it is the product, over the two colliding
indices, of a single universal function of one variable.  In the language of
\S\ref{sec:weight}: the $1$-cochain $\{(\bM,ab)\mapsto\rho\}$ on the collision
groupoid is the coboundary of the $0$-cochain $\Psi^{-1}$.  Had $\rho$ depended on
$w_c$ for some $c\notin\{a,b\}$, or failed to factor by variable, \emph{no}
renormalisation of this shape could trivialise it, and everything below would
fail.  Note also $\psi|_{q=0}=1$, so the renormalisation is invisible at $q=0$: it
is a purely $q$-theoretic twist, consistent with Rick's observation that
``at $q=0$ the twist is 1 (GKM).''
\end{rem}

%=====================================================================
\section{Collision weight, and the lattice lemma}
\label{sec:weight}

\begin{defn}\label{def:weight}
Fix $a<b$ and $d\in\mathbb{Z}$.  A function $f\in\mathcal{O}_{ab}^{\mathcal{P}}$
has \textbf{collision weight $d$ at $H_{ab}$} if
$f(\bM)|_{H}=g^{d}f(\bM'')|_{H}$ for every $\bM\supseteq\{a,b\}$.  Write
$W_d(ab)$ for the set of such $f$.
\end{defn}

\noindent
Immediately: each $W_d(ab)$ is an $\mathcal{O}_{ab}$-module (the scalar does not
depend on $\bM$); $W_d\cdot W_e\subseteq W_{d+e}$ under pointwise product; and
$W_0$ is a subring containing $1$.  By (QC), $\phi_{\bI}\in W_1(ab)$ for every
$\bI$; note in particular $1\notin W_1$ unless $g=1$, so $W_1$ is \emph{not} a
ring --- and that failure is the whole phenomenon.

\begin{prop}[lattice lemma; Rick, Prop 5, re-derived]\label{prop:lattice}
$\operatorname{span}_{\mathcal{O}_{ab}}\{\phi_{\bI}:\bI\subseteq[N]\}=W_1(ab)$.
Consequently, for $f\in\mathcal{O}_{ab}^{\mathcal{P}}$ with coordinates
$\ell_\nu(f)$ in the basis $\{\phi_{\bI}\}$:
\[
  \ell_\nu(f)\in\mathcal{O}_{ab}\ \text{for all }\nu
  \iff f\in W_1(ab),
\]
and any pole that does occur is simple.
\end{prop}

\begin{proof}
``$\subseteq$'': the $\phi_{\bI}$ are Laurent polynomials in $w$ over
$\mathbb{C}[q]$ localised at $(q)$, hence lie in
$\mathcal{O}_{ab}^{\mathcal{P}}$, and satisfy the congruence by (QC); $W_1$ is an
$\mathcal{O}_{ab}$-module.

Colength of $W_1$.  The \emph{defect map}
$D:\mathcal{O}_{ab}^{\mathcal{P}}\to k_{ab}^{\{\bM\supseteq\{a,b\}\}}$,
$D_{\bM}(f):=(f(\bM)-g\,f(\bM''))|_{H}$, is $\mathcal{O}_{ab}$-linear with kernel
$W_1$ by definition.  It is surjective: for $\bM_0\supseteq\{a,b\}$ take $f$ the
indicator of the single point $\bM_0$; then
$D_{\bM}(f)=\delta_{\bM,\bM_0}$, because $\bM''=\bM_0$ would force
$\bM=\bM_0\cup\{a,b\}=\bM_0$, excluded.  Hence
$\operatorname{length}(\mathcal{O}_{ab}^{\mathcal{P}}/W_1)=\#\{\bM\supseteq\{a,b\}\}=2^{N-2}$.

Colength of the span.  Over the DVR $\mathcal{O}_{ab}$ the colength of the
lattice spanned by the columns of $(\phi_{\bI}(\bM))_{\bM,\bI}$ is
$v_{ab}(\det)$.  By (M3) the matrix is triangular for Motegi's order, so
$\det=\pm\prod_{\bI}\At_{\bI,\bI}$, and by (M4)
$v_{ab}(\At_{\bI,\bI})=1$ if $a,b\in\bI$ (the single factor
$(1-w_a^{-1}w_b^{-1})$) and $0$ otherwise: at generic $q$ no other factor of the
positive-root product vanishes identically on $H_{ab}$ --- in particular
$(1-qw_aw_b)|_{H}=1-q\ne0$ and $(w_a-w_b)|_{H}=w_a-w_a^{-1}\not\equiv0$.  Hence
$v_{ab}(\det)=\#\{\bI\supseteq\{a,b\}\}=2^{N-2}$.

Two $\mathcal{O}_{ab}$-submodules of $\mathcal{O}_{ab}^{\mathcal{P}}$, one
contained in the other, with equal finite colength, are equal.  The displayed
equivalence follows: if all $\ell_\nu(f)$ are regular then
$f=\sum\ell_\nu(f)\phi_\nu$ lies in the span; conversely the span's elements have
regular coordinates by uniqueness.  Simplicity of the poles: $t_{ab}\cdot
\mathcal{O}_{ab}^{\mathcal{P}}\subseteq W_1$.
\end{proof}

\begin{rem}[the two numerical inputs, checked independently]
Both inputs to the colength count were verified directly against the engine's
printed diagonals at $N=3$, type $C$ (file
\texttt{code-1016-q492/q492\_N3\_census.log}).  For $H_{12}$ the factor
$(w_1w_2-1)$ occurs in exactly $\At_{\bI,\bI}$ for $\bI=(1,2)$ and
$\bI=(1,2,3)$, each once: $2=2^{3-2}$.  Likewise for $H_{13}$ and $H_{23}$.
\end{rem}

%=====================================================================
\section{Theorem B: the long-root poles are removable}
\label{sec:B}

\begin{thm}[Rick's hunch, confirmed for long roots]\label{thm:B}
Let $\chi\in\mathbb{C}(q)(v)^\times$ satisfy
\[
  \textbf{(COB)}\qquad \chi(x)\,\chi(1/x)=g(x)^{-1}
  \quad\text{identically in }x,
\]
let $\Psi_\chi(\bM):=\prod_{m\in\bM}\chi(w_m)$ and
$\widehat\phi_{\bI}:=\Psi_\chi\phi_{\bI}$, and assume
$\Psi_\chi\in\mathcal{O}_{ab}^{\mathcal{P}}$.  Let
$\widehat c^{\,\nu}_{\lambda\mu}$ be the structure constants of
$\{\widehat\phi_{\bI}\}$.  Then for every $\lambda,\mu,\nu$ and every $a<b$,
$\widehat c^{\,\nu}_{\lambda\mu}$ is regular at $H_{ab}$.

In particular $\chi=\psi$ satisfies \textup{(COB)} by Theorem~\ref{thm:A}, for
every pair $a<b$ at once.
\end{thm}

\begin{proof}
$\Psi:=\Psi_\chi$ is invertible in $\mathcal{A}$ (each $\Psi(\bM)\in F^\times$),
so
\[
  \widehat\phi_\lambda\widehat\phi_\mu=\sum_\nu \widehat c^{\,\nu}\widehat\phi_\nu
  \iff
  \Psi^2\phi_\lambda\phi_\mu=\Psi\sum_\nu \widehat c^{\,\nu}\phi_\nu
  \iff
  \Psi\cdot\phi_\lambda\phi_\mu=\sum_\nu \widehat c^{\,\nu}\phi_\nu,
\]
i.e.\ $\widehat c^{\,\nu}=\ell_\nu(f)$ with $f:=\Psi\cdot\phi_\lambda\phi_\mu$.
Now compute the weight of $f$ at $H_{ab}$.  Since
$\bM=\bM''\sqcup\{a,b\}$ and $w_b=x^{-1}$ on $H$,
\[
  \Psi(\bM)\big|_{H}=\chi(x)\chi(1/x)\,\Psi(\bM'')\big|_{H}=g^{-1}\Psi(\bM'')\big|_{H},
\]
so $\Psi\in W_{-1}(ab)$; and $\phi_\lambda\phi_\mu\in W_1\cdot W_1=W_2(ab)$ by
(QC).  Hence $f\in W_{-1}\cdot W_2=W_1(ab)$.  By
Proposition~\ref{prop:lattice}, every $\ell_\nu(f)$ is regular at $H_{ab}$.
\end{proof}

\begin{rem}[why this is not magic]
The mechanism is entirely visible in Rick's Theorem~6.  A product of two basis
elements collides with multiplier $g^2$ while the span collides with $g$; the
mismatch $g^2-g=g(g-1)$ \emph{is} the pole.  Multiplying everything by a
function of weight $-1$ moves the product to weight $1$ and the span to weight
$0$ --- no, more precisely: it moves the product $\Psi\phi_\lambda\phi_\mu$ to
weight $1$, which is the weight of the span, and the mismatch disappears.
Equivalently: $\{\widehat\phi_{\bI}\}$ spans a \emph{ring} of weight-$0$
functions, and a ring has regular structure constants.
\end{rem}

%=====================================================================
\section{The short-root condition, which the data forced}
\label{sec:short}

Theorem~\ref{thm:B} says nothing about the other divisors of $\det\At$.  I did not
expect that to matter, and it does.  The $\chi=v\,\psi(v)$ arm --- chosen because
$\eta(v)=v$ satisfies $\eta(x)\eta(1/x)=1$ and therefore preserves (COB) --- removed
\emph{every} long-root pole as predicted and produced $57$ new ones, all on
divisors $\{w_i=-1\}$.

\begin{thm}[(SRC) short-root collapse, type $C$]\label{thm:SRC}
For every $N$, every $\bI$, every $\bM\ni m$ and every $\varepsilon$ with
$\varepsilon^2=1$,
\[
  \At_{\bM,\bI}\big|_{w_m=\varepsilon}=\At_{\bM^{-},\bI}\big|_{w_m=\varepsilon}:
\]
the short-root collision multiplier is exactly~$1$, uniformly in $\bM$ and $\bI$.
\end{thm}

\begin{proof}
Write $s:=|\bM|$, $r:=|\bI|$, $R:=\bM^-$, and run Motegi's symmetrisation formula
unpadded at $n=s$ and at $n=s-1$ (stability, L2823--2848), with $\beta=1$,
$\gamma=0$.  Throughout, $\varepsilon^{-1}=\varepsilon$.

\emph{Step 1: $\varepsilon$ must be a tail.}  The head factor of line~1 is
$(1-q)(u_j^2-1)$, which vanishes at $u_j=\varepsilon$.  So every permutation
placing $\varepsilon$ in a head position contributes $0$.  The $n-r$ tail
positions are symmetric, so the $n-r$ ways of seating $\varepsilon$ in a tail
give equal contributions and cancel the change $1/(n-r)!\to1/(n-1-r)!$.  The
surviving head selections are exactly the head selections of the $(s-1)$-formula.
(If $s-1<r$ there is no tail slot at all, every permutation dies, and both sides
are $0$.)

\emph{Step 2: the extra factors are head-independent.}  Fix a head selection and
compare the $s$-term with the $(s-1)$-term.  The $w$-alphabet is unchanged, so
line~1 matches.  The extra factors carried by the tail $\varepsilon$ are
\[
  \underbrace{\textstyle\prod_{k=1}^{N}(\varepsilon-qw_k)}_{\text{line 2b}}\cdot
  \underbrace{\textstyle\prod_{\text{heads }h}\frac{q\varepsilon-h}{\varepsilon-h}(\varepsilon h-1)}_{\text{line 3a}}\cdot
  \underbrace{\textstyle\prod_{\text{tails }t\ne\varepsilon}(\varepsilon t-q)}_{\text{line 3b}}\cdot
  \underbrace{\varepsilon}_{\text{line 3c}} .
\]
Now $\varepsilon h-1=\varepsilon(h-\varepsilon)=-\varepsilon(\varepsilon-h)$, so the
line-3a factor collapses to $(-\varepsilon)^{r}\prod_h(q\varepsilon-h)
=\varepsilon^{r}\prod_h(h-q\varepsilon)$; and
$\varepsilon t-q=\varepsilon(t-q\varepsilon)$, so line 3b is
$\varepsilon^{\,n-r-1}\prod_{t\ne\varepsilon}(t-q\varepsilon)$.  The head product
and the tail product therefore \textbf{recombine into a product over all of $R$}:
\[
  \varepsilon^{r}\cdot\varepsilon^{\,n-r-1}\prod_{v\in R}(v-q\varepsilon)
  =\varepsilon^{\,s-1}\prod_{v\in R}(v-q\varepsilon),
\]
with no trace of the head selection.  This is the exact analogue of step~2 of
Rick's Lemma QC, and it is the only step that could have failed.

\emph{Step 3: the ratio.}  The prefactor $\prod_k u_k^{-(N+n)}$ contributes
$\varepsilon^{-(N+s)}\prod_{v\in R}v^{-1}$.  Collecting, and using
$v^{-1}(v-q\varepsilon)=1-q\varepsilon/v$ and
$\varepsilon^{-N}\prod_k(\varepsilon-qw_k)=\prod_k(1-q\varepsilon w_k)$,
\[
  \frac{F_{s}}{F_{s-1}}\bigg|_{w_m=\varepsilon}
  =\prod_{k=1}^{N}(1-q\varepsilon w_k)\ \prod_{v\in R}\Bigl(1-\frac{q\varepsilon}{v}\Bigr).
  \tag{$\ast$}
\]

\emph{Step 4: the normaliser does the same.}  Adjoining $u=\varepsilon$ to
$\mathcal{E}^C$ adds, to the numerator,
$\prod_{v\in R}(1-qv/\varepsilon)(1-q\varepsilon/v)\cdot(1-q)$ from
$\prod_{j,k}(1-qu_k/u_j)$ and $\prod_k(1-q\varepsilon w_k)$ from
$\prod_{j,k}(1-qu_jw_k)$; and, to the denominator,
$\prod_{v\in R}(1-q\varepsilon v)$ from $\prod_{j<k}(1-qu_ju_k)$ together with one
factor $(1-q)$ from $(1-q)^{n-r}$.  The two $(1-q)$'s cancel, and
$(1-qv/\varepsilon)=(1-q\varepsilon v)$ cancels against the $j<k$ product, leaving
\[
  \frac{\mathcal{E}^C_{s}}{\mathcal{E}^C_{s-1}}\bigg|_{w_m=\varepsilon}
  =\prod_{k=1}^{N}(1-q\varepsilon w_k)\ \prod_{v\in R}\Bigl(1-\frac{q\varepsilon}{v}\Bigr),
\]
which is $(\ast)$.  Hence $\At_s/\At_{s-1}=1$.
\end{proof}

\begin{rem}[verification, with a plant]
The two ratios of Steps 3 and 4 were checked \emph{separately} against their
closed forms --- not merely their quotient, which could be $1$ with both wrong ---
over all $(N,\bM,m,\bI,\varepsilon)$ with $N\le3$ and $s-1\ge r$: $112/112$ for
$F$, $112/112$ for $\mathcal{E}^C$, $112/112$ for $\At$ itself
(\texttt{src\_proof\_check.py}).  The proof uses $\varepsilon^2=1$ in three
places, so the plant is $\varepsilon=2$: the $\At$ equality must \emph{fail}
there, and it fails $24/24$.  A green at $\varepsilon=2$ would have meant the
instrument was blind.
\end{rem}

With multiplier $1$ the weight filtration collapses ($W_d=W_0$ for all $d$), so
$\phi_\lambda\phi_\mu$ and the span sit at the same weight and the baseline
constants have no short-root poles --- which is what the baseline census shows.
A renormalisation has short-root weight $\chi(\varepsilon)$:
\[
  \Psi_\chi(\bM)\big|_{w_m=\varepsilon}
  =\chi(\varepsilon)\,\Psi_\chi(\bM^{-})\big|_{w_m=\varepsilon},
\]
so the condition for $\Psi_\chi\phi_\lambda\phi_\mu$ to stay in the short-root
lattice is
\[
  \textbf{(SR)}\qquad \chi(\varepsilon)=1 \quad\text{for every }\varepsilon
  \text{ with }\varepsilon^{h}=1 .
\]

\begin{prop}[short-root lattice lemma]\label{prop:short}
$\operatorname{span}_{\mathcal{O}_{m,\varepsilon}}\{\phi_{\bI}\}
 =\{f:(f(\bM)-f(\bM^{-}))|_{S_m^\varepsilon}=0\ \forall \bM\ni m\}$,
assuming \textup{(SRC)} = Theorem~\ref{thm:SRC}.  Consequently if $\chi$ satisfies \textup{(SR)} then
every $\widehat c^{\,\nu}_{\lambda\mu}$ is regular at $S_m^{\varepsilon}$.
\end{prop}
\begin{proof}
Verbatim the proof of Proposition~\ref{prop:lattice} with $g$ replaced by $1$ and
$\bM''$ by $\bM^-$.  The two counts: the defect map onto
$k^{\{\bM\ni m\}}$ is surjective by the same indicator argument, giving colength
$2^{N-1}$; and $v_{S_m^\varepsilon}(\At_{\bI,\bI})=1$ iff $m\in\bI$ --- the
factor $(1-w_m^{-2})$, which is the $j=k$ term of (M4)'s $\prod_{j\le k}$ --- and
no other positive-root factor vanishes on $\{w_m=\varepsilon\}$ at generic $q$,
since $(1-qw_j^2)$ needs $w_j^2=q^{-1}$ and the $A$-type and $q$-pair factors cut
out different divisors.  So $v(\det)=\#\{\bI\ni m\}=2^{N-1}$.  For the
consequence: $f=\Psi_\chi\phi_\lambda\phi_\mu$ has
$f(\bM)|_{S}=\chi(\varepsilon)\cdot1\cdot f(\bM^-)|_{S}=f(\bM^-)|_{S}$.
\end{proof}

\begin{rem}[inevitability]\label{rem:inev}
The two conditions are met by $\psi$ \emph{for the same reason}.  From
\texttt{eq:dict},
$\rho^X-1=q(t-1)^2/((1-qt)(1-q/t))$ with $t=w_a^h$, so
$\{\rho^X=1\}=\{q=0\}\cup\{w_a^{h}=1\}$: the obstruction vanishes exactly on the
$h$-th roots of unity.  And $\psi(v)=(1-qv^h)/(1-q)$ satisfies $\psi(\varepsilon)=1$
exactly when $\varepsilon^{h}=1$.  \textbf{The locus where the multiplier is
already trivial is precisely the locus where the renormaliser is required to be
trivial.}  The short-root divisors of $\det\At$ are $\{w_m^{h}=1\}$ --- the same
set.  I did not arrange this; I ran an arm that broke, and the diagnosis
predicted that it must break at $\varepsilon=-1$ and \emph{not} at
$\varepsilon=+1$, which is what the data show.
\end{rem}

%=====================================================================
\section{Theorem C: the exact family}

\begin{thm}[exact characterisation]\label{thm:C}
Let $\chi\in\mathbb{C}(q)(v)^\times$ be regular and non-vanishing at the generic
point of each divisor below, $\Psi=\Psi_\chi$, and put
$\kappa:=\chi(x)\chi(1/x)\in k_{ab}^\times$.  Fix $\lambda,\mu$.
\begin{enumerate}
\item \textup{(long roots)}  For $a<b$, the following are equivalent:
 \begin{enumerate}
  \item every $\widehat c^{\,\nu}_{\lambda\mu}$ is regular at $H_{ab}$;
  \item $(\kappa g-1)\cdot(\phi_\lambda\phi_\mu)(\bM'')|_{H}=0$ for every
        $\bM''\subseteq[N]\sm\{a,b\}$;
  \item $\kappa=g^{-1}$, \emph{or} no $\bM''\subseteq[N]\sm\{a,b\}$ satisfies both
        $\bM''\ge\lambda$ and $\bM''\ge\mu$.
 \end{enumerate}
\item \textup{(short roots)}  For $\varepsilon^h=1$, every
  $\widehat c^{\,\nu}_{\lambda\mu}$ is regular at $S_m^{\varepsilon}$ if
  $\chi(\varepsilon)=1$; and if $\chi(\varepsilon)\ne1$ it is regular there only
  when $(\phi_\lambda\phi_\mu)(\bM^-)|_{S}=0$ for all $\bM\ni m$.
\end{enumerate}
Hence the renormalisations that are pole-free on \emph{all} these divisors for
\emph{all} $(\lambda,\mu)$ are exactly
\[
  \boxed{\ \chi=\psi\cdot\eta,\qquad
   \eta(x)\eta(1/x)=1,\qquad \eta(\varepsilon)=1\ \ \forall\,\varepsilon^{h}=1.\ }
\]
\end{thm}

\begin{proof}
(1) Put $f:=\Psi\phi_\lambda\phi_\mu$.  By (QC) and the computation in
Theorem~\ref{thm:B}, $f(\bM)|_{H}=\kappa g^2 f(\bM'')|_{H}$, so the defect is
$D_{\bM}(f)=(\kappa g^2-g)f(\bM'')|_{H}=g(\kappa g-1)f(\bM'')|_{H}$.  By
Proposition~\ref{prop:lattice}, (a) $\iff$ $D_{\bM}(f)=0$ for all $\bM$, and
$g\ne0$ gives (b).  For (b)$\iff$(c): if $\kappa=g^{-1}$ then $\kappa g-1=0$.
Otherwise $\kappa g-1\ne0$ in $k_{ab}$, and since $\Psi(\bM'')|_H\ne0$ by
hypothesis we need $(\phi_\lambda\phi_\mu)(\bM'')|_{H}=0$ for all $\bM''$.  That
holds iff no $\bM''\ge\lambda,\mu$: ``$\Leftarrow$'' is (M3); ``$\Rightarrow$'' is
Rick's Lemma L (his \S2bis: a degeneration $w_c=t^{-e_c}$, $t\to0^+$, in which
$\phi^0_\kappa(\bM'')\to1$ when $\bM''\ge\kappa$), which shows
$\phi^0_\lambda\phi^0_\mu(\bM'')|_{H}\ne0$ at $q=0$ and hence
$\ne0$ at generic $q$.

(2) is Theorem~\ref{thm:SRC} and Proposition~\ref{prop:short}, with the same defect computation at $g=1$.

For the final display: (COB) is $\kappa=g^{-1}$, and $\psi$ is one solution by
Theorem~\ref{thm:A}, so the solution set of (COB) is the coset
$\psi\cdot\{\eta:\eta(x)\eta(1/x)=1\}$; (SR) is $\chi(\varepsilon)=1$, and
$\psi(\varepsilon)=1$ for $\varepsilon^h=1$ (Remark~\ref{rem:inev}), so it
transfers to $\eta(\varepsilon)=1$.
\end{proof}

\begin{cor}[the brief's Theorem C, answered negatively]
The long-root pole order is \textbf{not} invariant across the family of diagonal
renormalisations.  By Proposition~\ref{prop:lattice} it is at most $1$, and by
Theorem~\ref{thm:C}(1) it is $0$ on the coset $\psi\cdot\{\eta\}$ and $1$ off it
(on every non-degenerate $(\lambda,\mu)$).  The brief asked for an invariance
theorem that would bind Rick's unwritten geometric normalisation; instead the
pole-removing normalisations are pinned down exactly, and any geometric
normalisation that is diagonal in this sense \emph{must} lie in the boxed coset.
\end{cor}

\begin{rem}[monomial twists]
$\eta(v)=v^{k}$ satisfies $\eta(x)\eta(1/x)=1$ for every $k$, and
$\eta(\varepsilon)=\varepsilon^{k}$, so in type $C$ ($h=2$,
$\varepsilon=\pm1$) the admissible monomial twists are exactly the \emph{even}
powers.  This was used as a falsifiable prediction: $\chi=v\psi(v)$ must fail at
$w=-1$ only, and $\chi=v^2\psi(v)$ must be clean.  See \S\ref{sec:verif}.
\end{rem}

%=====================================================================
\section{Verification}
\label{sec:verif}

All computations: pure SymPy, engine
\texttt{proofs/code-1017c2-coboundary/q484\_engine.py} (the 2026-10-17 copy, the
only one carrying $\mathcal{E}^B$; for type $C$ it is byte-identical in the
relevant code path to the copies used on 10-15 and 10-16, checked by \texttt{diff}).
Census: every $(\lambda,\mu,\nu)$ with $\mu\ge\lambda$ in the linear extension
and $c^\nu_{\lambda\mu}\ne0$; $N=3$, type $C$; $115$ non-zero constants.
A constant is counted \emph{non-Laurent} iff, after \texttt{cancel}, its
denominator has a non-monomial irreducible factor.

\paragraph{Instrument controls.}
\begin{itemize}
\item \textbf{Baseline reproduces.}  $\chi=1$ at symbolic $q$ returns
  $13/115$ non-Laurent, which is the 2026-10-16 census number, and the poles are
  on long-root divisors only.
\item \textbf{The $q$-specialisation is faithful.}  $\chi=1$ at $q=3/7$ returns
  $13/115$ \emph{and the symmetric difference against the symbolic run is
  empty} --- the same $13$ triples with the same denominators, not merely the
  same count.  This licenses reading the other arms at $q=3/7$.
\item \textbf{The instrument can see movement in both directions}: one arm
  removes poles ($\psi$), one arm adds them ($v\psi$).  A green that could only
  ever be green is not evidence.
\end{itemize}

\paragraph{Arms.}  Each $\chi$ gives $\Psi_\chi(\bM)=\prod_{m\in\bM}\chi(w_m)$ and
a full recomputation of all $115$ constants from $\Ah$ by the triangular solve
--- an independent route to $\widehat c$, not the $\ell_\nu(\Psi\phi_\lambda\phi_\mu)$
rewriting used in the proof.  Predictions were written to
\texttt{PREDICTIONS.md} before the arms finished.

\begin{center}
\begin{tabular}{@{}llccl@{}}
\toprule
arm & $\chi(v)$ & (COB) & (SR) & result \\
\midrule
\texttt{id}    & $1$                    & fail  & holds       & $13/115$; long-root only \\
\texttt{psi}   & $(1-qv^2)/(1-q)$       & holds & holds       & $\mathbf{0/115}$ \\
\texttt{x2psi} & $v^2(1-qv^2)/(1-q)$    & holds & holds       & $\mathbf{0/115}$ \\
\texttt{xpsi}  & $v(1-qv^2)/(1-q)$      & holds & fails at $-1$ & $57/115$; \textbf{zero} long-root, all at $w_i+1$ \\
\texttt{x}     & $v$                    & fail  & fails at $-1$ & $57/115$; long-root \emph{and} $w_i+1$ \\
\texttt{sq}    & $((1-qv^2)/(1-q))^2$   & fail  & holds       & $13/115$; \textbf{the same 13, same denominators} \\
\texttt{h1}    & $(1-qv)/(1-q)$         & fail  & fails at $-1$ & $57/115$; long-root \emph{and} $w_i+1$ \\
\texttt{unrm}  & $1-qv^2$               & fail  & fails at $\pm1$ & $57/115$; long-root, $w_i+1$ \emph{and} $w_i-1$ \\
\bottomrule
\end{tabular}
\end{center}

\paragraph{Every pre-registered prediction holds, at divisor level.}
The table's ``result'' column is not a count: each entry records \emph{which
divisor classes} carry poles.  Reading the two condition columns against the
three divisor classes:
\begin{itemize}
\item \textbf{(COB) governs the long-root divisors and nothing else.}  The four
  arms with (COB) failing (\texttt{id}, \texttt{x}, \texttt{sq}, \texttt{h1},
  \texttt{unrm}) all have long-root poles; the three with (COB) holding
  (\texttt{psi}, \texttt{x2psi}, \texttt{xpsi}) have \emph{none}, including
  \texttt{xpsi}, which is riddled with poles of the other kind.
\item \textbf{(SR) governs the short-root divisors, and separately at each
  $\varepsilon$.}  The arms with $\chi(-1)\ne1$ (\texttt{xpsi}, \texttt{x},
  \texttt{h1}, \texttt{unrm}) have poles at $w_i+1$ and no others have any;
  \texttt{unrm} is the \emph{only} arm with $\chi(+1)\ne1$ and the \emph{only}
  arm, of eight, that produces a $w_i-1$ pole.
\item \textbf{The positive prediction was made before it was measured.}
  $\chi=v^2\psi$ satisfies both conditions and was predicted clean at a value of
  $\chi$ not previously run; it returns $0/115$.
\item \textbf{The sharpest set prediction holds.}  Theorem~\ref{thm:C}(1) makes
  the long-root pole locus depend on $\chi$ only through the predicate
  $\kappa g=1$, which fails for both \texttt{id} and \texttt{sq}; so \texttt{sq}
  must return not merely $13$ but \emph{the same} $13$ triples.  It does, with
  identical denominators on every one.
\end{itemize}

\paragraph{A grading lesson, paid for in this session.}
\texttt{x} and \texttt{xpsi} both return $57$, and their sets of $57$ triples are
\textbf{identical} --- symmetric difference empty.  Grading by count would have
called them the same; grading by set, which is my standing rule
(\texttt{two-criteria-can-fit-the-same-count-and-differ-as-sets}), would
\emph{also} have called them the same.  They are not: on exactly the $13$
baseline triples their denominators differ, \texttt{x} retaining the long-root
factor that \texttt{xpsi} has shed, e.g.\ at $(\emptyset,\emptyset,\{1,2\})$
\[
  \texttt{x}:\ (w_2+1)(w_1+1)(w_1w_2-1)
  \qquad\text{versus}\qquad
  \texttt{xpsi}:\ (w_2+1)(w_1+1).
\]
The whole of Theorem~\ref{thm:B} lives in that difference.  So: \emph{a set is
itself a projection.  When two arms agree as sets, descend to the label each
element carries} --- here, the pole divisor, not the triple.  One level below the
count-to-set rule, and the first time it has mattered.

\paragraph{Generic $q$.}  The main positive arm was also run at symbolic $q$
rather than $q=3/7$: $\chi=\psi$ over
$\mathbb{C}(q,w_1,w_2,w_3)$ returns $\mathbf{0/115}$ as well, so
Theorem~\ref{thm:B} is confirmed at generic $q$ and not only at a specialisation.
The baseline was likewise run at symbolic $q$ ($13/115$).  The other six arms were
read at $q=3/7$, licensed by the empty symmetric difference of the two baseline
runs.


%=====================================================================
\section{Where the pre-registered mechanism broke}
\label{sec:error}

The brief's Theorem-B argument reads $\prod_{m\in\bM}\psi(w_m)$ as a
\emph{diagonal rescaling of the basis}, and concludes that a structure constant
is multiplied by a ratio of $\psi$'s --- a unit at $H_{ab}$, hence unable to
cancel a pole.  That conclusion is correct for the operation it names, and the
operation it names is not the one in the formula.

\begin{itemize}
\item \textbf{Label-indexed} rescaling: $\phi_{\bI}\mapsto\chi_{\bI}\phi_{\bI}$
  with $\chi_{\bI}\in F$.  Then
  $\widehat c^{\,\nu}_{\lambda\mu}=(\chi_\lambda\chi_\mu/\chi_\nu)\,c^{\nu}_{\lambda\mu}$,
  a unit multiple; poles cannot move.  \emph{True.}
\item \textbf{Point-indexed} rescaling: $\phi_{\bI}\mapsto\Psi\phi_{\bI}$ with
  $\Psi\in\mathcal{A}$ a single function of the \emph{point}.  Then
  $\widehat c^{\,\nu}_{\lambda\mu}=\ell_\nu(\Psi\,\phi_\lambda\phi_\mu)$, which is
  not a unit multiple of $c^{\nu}_{\lambda\mu}$ at all: multiplication by a fixed
  element of the algebra reshuffles coordinates.
\end{itemize}

The product $\prod_{m\in\bM}$ is indexed by $\bM$, the point.  And the decisive
point is that \textbf{Theorem~A requires the point-indexed reading}: the collapse
$\bM\mapsto\bM''$ acts on the point, so a label-indexed $\chi_{\bI}$ cancels out
of the collapse ratio entirely and changes no multiplier.  So the brief proved
Theorem~A under reading~2 and refuted Theorem~B under reading~1, in adjacent
paragraphs, with the same symbol.

\begin{quote}
\textbf{Rule.}  The index set of a rescaling decides whether it is a change of
basis or multiplication by an algebra element.  Before reasoning about what a
rescaling does to structure constants, say out loud which set the product runs
over --- and check that the \emph{other} theorems in the same document are using
the same reading.
\end{quote}

This is a relative of the one-coordinate-match family
(\texttt{the-missing-coordinate-is-one-rule-in-nine-drawers}), but it is a
different animal and I want the distinction recorded.  There, two \emph{objects}
agreed on one coordinate and I inferred identity.  Here there is one object and
two \emph{readings of my own notation}, and the brief's two theorems required
\emph{incompatible} readings: the document was internally inconsistent while each
paragraph was locally correct.  A consistency check across paragraphs of one's
own brief would have caught it; a check of any single paragraph would not.

%=====================================================================
\section{Consequences}

\begin{retr}\label{retr:q491}
Q491's criterion --- $c^{\nu}_{\lambda\mu}$ Laurent $\iff$
$|\nu|-\max(|\lambda|,|\mu|)\le1$, measured $225/225$ with zero false positives
and zero false negatives --- is a property of the $\mathcal{E}^C$-normalised
basis $\At$, not of the algebra $\mathcal{A}$.  At $N=3$ every non-Laurent
constant has its poles on long-root divisors only, so under the $\psi$
renormalisation all $115$ become Laurent and the criterion becomes vacuous.  The
criterion remains true and useful \emph{of $\At$}; what is withdrawn is any
reading of it as intrinsic.
\end{retr}

\begin{retr}\label{retr:q500}
The negative half of Q500 --- ``pole $\iff$ Pfaffian category'', with the pole as
an obstruction --- is withdrawn in the form I stated it.  The pole is not an
invariant of the algebra; it is an invariant of a choice of normaliser, and the
choice that removes it is explicitly available.  The identity
$\rho^X-1=q\,w_a^h([K^X(w_a)]_0^1)^2$ is unaffected: it is a true statement about
the \emph{cochain} $\rho$.
\end{retr}

\noindent
What replaces them is better, and is the actual content of today:

\begin{quote}
The obstruction is a \textbf{cohomology class}, and it is \textbf{zero}.  The
poles are removable \emph{because} $\rho$ splits as $\varphi(w_a)\varphi(w_b)$.
Had it not split --- had it depended on a third $w_c$, or failed to factor by
variable --- Theorem~\ref{thm:C} says no point-diagonal renormalisation could
have removed the poles, and the obstruction would have been real.  The
non-vanishing of $\rho-1$ and the existence of the poles are both statements
about a representing cochain.  What I had been calling the obstruction was the
cochain.
\end{quote}

\noindent
This also sharpens Rick's sentence.  He wrote that my poles \emph{measure the
normalisation defect}.  Strictly that identifies two objects --- the defect of
the \emph{basis} ($\rho\ne1$) and the poles of the \emph{structure constants} ---
and my pre-registered objection to the identification was the right instinct
applied to the wrong case: they are different objects, and here they happen to be
controlled by the same cochain, so trivialising one trivialises the other.  The
correct statement is the implication, not the identity: \emph{$\rho$ a coboundary
$\Rightarrow$ the structure-constant poles are removable}, with Theorem~\ref{thm:C}
supplying the converse.

%=====================================================================
\section{Gaps and questions}

\begin{gapthm}[type $B$, and $N\ge4$]
Theorems~\ref{thm:A}--\ref{thm:C} are stated for general $h$ and use only (QC),
(M3), (M4) and (SRC); the type $B$ input (QC) is proved
(\texttt{thm:rhoB}).  But the \emph{census} is type $C$, $N=3$ only.  Q502 (re-measure
every type $B$ number against $\mathcal{E}^B$) is untouched.  At $N=4$ Rick
reports that the first-order route behaves differently (his (b) and (c) both fail
first at $N=4$); nothing in \S\ref{sec:B} is first-order, so I expect
Theorem~\ref{thm:B} to survive, but it is unmeasured.
\end{gapthm}

\begin{gapthm}[other divisors]
$\det\At$ also vanishes on the $A$-type divisors $\{w_i=w_j\}$ and on
$q$-divisors $\{q\cdot\text{monomial}=1\}$.  No baseline constant has a pole
there, and no arm produced one, but I have no theorem.
\end{gapthm}

\paragraph{New questions.}
\textbf{Q510} --- \emph{is $\Ah$ the weight function of a model?}  $\Ah$ is a
basis with honest GKM collapse at both long and short roots and a genuine
multiplication rule with no poles on them.  Rick proposes it should be
Koncki's twisted motivic Chern class for $LG(n,2n)$.  The test is cheap and
sharp: $\psi|_{q=0}=1$, so $\Ah|_{q=0}=\At|_{q=0}$, and the two can only differ
in their $q$-deformation; compare $\Ah$ against Koncki's table for $LG(2,4)$ on
the entry Rick predicts is the odd one out, $(\{1,2\},\emptyset)$.  He states he
has not checked Koncki's table; neither have I.  Held at
\texttt{agent-summary}; a first-hand read is the prerequisite.

\medskip\noindent
\textbf{Gate Q504 observed.}  No novelty sentence is made here about Q484, Q491
or Q492.  Digalaki--Wheeler, \emph{Structure Constants for Free-Fermion Six
Vertex Model Symmetric Functions} (Mittag-Leffler, July 2026, in preparation, no
arXiv ID) is held only at \texttt{agent-summary}, and their hypothesis is the
free-fermion locus, which I conjecture is my $\rho=1$.  Theorems~\ref{thm:A},
\ref{thm:B} and \ref{thm:C} concern my own normaliser and the removability of my
own poles; Retractions~\ref{retr:q491} and \ref{retr:q500} are retractions of my
own claims.  Neither is a priority claim.

%=====================================================================
\section*{Acknowledgement}

Lemma QC, Proposition~5 and Theorem~6 of
\texttt{2026-10-17-q488-converse-order-q1-proofs.pdf} are Rick's.  The lattice
lemma above is his Proposition~5, re-derived here with its two numerical inputs
checked independently against my engine; Theorem~\ref{thm:C}(1)(c) uses his
Lemma~L.  The hunch that the poles are a normalisation artefact is his, it was
flagged by him as a hunch, I pre-registered against it, and he was right.

\end{document}
